\chapter{Conclusions}
In this laboratory work, I implemented the last two steps of the DES
round function and update needed to find the new right half. The
assignment starts with the combined output of the eight S-boxes, so
the same calculation can be used for any round once that word and
the previous left half are available.

The permutation required careful attention to bit numbering. The DES
table counts positions from the most significant bit, while the C\#
code uses right shifts to read a bit. Translating a table position
$p$ into a shift of $32-p$ lets the program build the output word in
the order specified by the table.

The XOR then combines the permuted word with the previous left half.
Its reversibility is part of the Feistel structure. Applying the
same XOR to $R_k$ recovers $L_{k-1}$, and the previous right half is
available as $L_k$. This connects the two implemented steps to the
way DES reverses a round.

The random mode provides inputs for a quick demonstration. To repeat
that calculation, I can enter the two displayed hexadecimal values
in manual mode. Showing the binary and hexadecimal forms together
makes it easier to follow the bit operations and read the final
32-bit result.
